\subsection{Ampliación de datos y gobernanza de la información}

La versión actual de ViaData — Medellín opera sobre un periodo Mede aproximado de 2014--2021 y un flujo ETL \textbf{manual pero reproducible}. Como línea prioritaria de trabajo futuro se plantea la \textbf{ingesta automática o programada} de nuevas versiones del conjunto Mede, con validación de esquema, alertas de calidad y versionado de cargas. Complementariamente, la \textbf{extensión temporal} del repositorio permitiría recalibrar modelos predictivos, reevaluar hold-out y contrastar tendencias post-2021, condicionado a la disponibilidad y consistencia de datos abiertos actualizados por la institución \cite{atancuri2024,oviedo2025}.

En paralelo, conviene fortalecer la gobernanza del dato: trazabilidad de transformaciones, metadatos de cada carga y documentación de cambios en catálogos o definiciones operativas, de modo que las proyecciones futuras mantengan comparabilidad con las evaluaciones ya cerradas.

\subsection{Indicadores y análisis geoespacial omitidos}

Durante el cronograma del proyecto se omitieron deliberadamente funcionalidades ligadas a catálogos que Mede no alimenta en la versión cargada:

\begin{itemize}
	\item \textbf{P11 — ranking por vía o punto crítico}, por ausencia de registros en tablas \texttt{via} y \texttt{punto\_critico}.
	\item \textbf{G04 — análisis por buffer de punto crítico}, mientras la tabla de puntos críticos permanezca vacía.
	\item \textbf{G05 — filtro por bounding box en mapa}, considerado redundante frente a filtros por comuna y barrio.
\end{itemize}

Una extensión natural del sistema consiste en incorporar estos indicadores \textbf{si la fuente institucional publica geometrías o identificadores de corredores e intersecciones} de forma estructurada, lo que habilitaría focalizar intervenciones en glorietas, tramos viales o nodos de alta siniestralidad ---línea de investigación coherente con estudios de puntos críticos en Medellín \cite{viveros2022,xu2022}.

Asimismo, profundizar el contraste entre modo \textit{registro} y modo \textit{espacial} (indicador G03) con reportes periódicos de calidad cartográfica aportaría evidencia para mejorar la georreferenciación en la captura administrativa.

\subsection{Modelos predictivos y evaluación}

La evaluación cerrada priorizó modelos \textbf{transparentes y comparables} (OLS, estacional, Poisson, media móvil, $\mu \pm 3\sigma$, ARIMA/SARIMA) y demostró que el desempeño depende del escenario analítico. Trabajo futuro en esta línea podría incluir:

\begin{itemize}
	\item \textbf{Modelos jerárquicos o bayesianos} para series territoriales con pocos datos por barrio, mitigando la limitación actual en que 220 de 270 barrios carecen de proyección individual en §3.
	\item \textbf{P15 — aprendizaje automático espacial} (por ejemplo, gradient boosting o modelos de conteo espacio-temporales), evaluados siempre con hold-out y sin sustituir la trazabilidad de las formulaciones actuales.
	\item \textbf{Variables exógenas} (meteorología, intervenciones viales, cambios normativos) para acercar las proyecciones a escenarios contrafactuales, sin pretender inferencia causal plena desde el registro administrativo.
	\item \textbf{Reevaluación periódica} del ranking de modelos al incorporar nuevos meses, automatizando los scripts \texttt{llenar\_evaluacion\_seccion*.py} en un pipeline de CI.
\end{itemize}

Se mantiene como criterio que cualquier modelo nuevo deba superar o complementar el protocolo MAPE hold-out $\leq 20\,\%$ antes de adoptarse como predeterminado en producción.

\subsection{Plataforma, despliegue y adopción institucional}

En el plano de ingeniería de software, extensiones viables incluyen:

\begin{itemize}
	\item \textbf{Generación de PDF en servidor} (Puppeteer, WeasyPrint u equivalente) como alternativa a la impresión nativa del navegador.
	\item \textbf{Activación del rol autoridad} reservado en el esquema de permisos, con flujos de aprobación o publicación de reportes oficiales.
	\item \textbf{Integración con tableros institucionales} o APIs de movilidad urbana, si la entidad dispone de canales de intercambio seguros.
	\item \textbf{Pruebas de carga y despliegue productivo} (contenedorización, respaldos, monitoreo) que exceden el alcance de validación con \texttt{pytest} del trabajo de grado.
\end{itemize}

El asistente conversacional basado en Gemini podría evolucionar hacia un catálogo ampliado de herramientas, auditoría de consultas y, en entornos institucionales, despliegue con políticas de privacidad y retención distintas al tier gratuito actual.

\subsection{Síntesis de prioridades}

En síntesis, las prioridades de trabajo futuro se ordenan así: \textbf{(1)} actualizar y automatizar el dato; \textbf{(2)} incorporar indicadores viales y puntos críticos si la fuente lo permite; \textbf{(3)} ampliar el repertorio predictivo con evaluación hold-out rigurosa, especialmente en barrios; y \textbf{(4)} madurar el despliegue y la adopción institucional del instrumento. La Figura~\ref{fig:flujo-trabajo-futuro} resume el flujo propuesto y las dependencias entre fases. Estas líneas prolongan el equilibrio alcanzado en ViaData — Medellín entre utilidad exploratoria, transparencia metodológica y reconocimiento explícito de limitaciones \cite{lugo2025,monar2025}.

\begin{figure}[H]
	\centering
	\resizebox{0.88\linewidth}{!}{%
	\begin{tikzpicture}[
		node distance=0.55cm,
		box/.style={draw, rounded corners, align=center, font=\scriptsize, inner sep=4pt, minimum width=6.2cm},
		dec/.style={draw, diamond, aspect=2.2, align=center, font=\scriptsize, inner sep=1pt, minimum width=2.8cm},
		end/.style={draw, rounded corners, align=center, font=\scriptsize, inner sep=4pt, minimum width=5.6cm, fill=gray!12},
		arr/.style={-{Stealth[length=1.8mm]}, thick},
		skip/.style={-{Stealth[length=1.8mm]}, thick, dashed}
	]
		\node[box, fill=blue!8] (act) {Estado actual: ViaData — Medellín (2014--2021, ETL manual reproducible)};
		\node[box, fill=green!12, below=of act] (p1) {\textbf{(1) Datos y gobernanza}\\Ingesta automática · validación de esquema · versionado · extensión temporal · metadatos};
		\node[dec, below=of p1] (cond) {¿Mede publica\\vías y puntos críticos?};
		\node[box, fill=yellow!18, below left=0.9cm and 0.2cm of cond] (p2) {\textbf{(2) Indicadores geoespaciales}\\P11 ranking vial · G04 buffer · G05 bbox · contraste registro/espacial};
		\node[box, fill=orange!15, below=1.35cm of cond] (p3) {\textbf{(3) Modelos predictivos}\\Jerárquicos/bayesianos · P15 ML espacial · variables exógenas · pipeline CI · MAPE hold-out $\leq 20\,\%$};
		\node[box, fill=violet!10, below=of p3] (p4) {\textbf{(4) Plataforma e institucionalización}\\PDF en servidor · rol autoridad · APIs · despliegue productivo · asistente ampliado};
		\node[end, below=of p4] (fin) {Sistema actualizado, trazable y orientado a adopción institucional};

		\draw[arr] (act) -- (p1);
		\draw[arr] (p1) -- (cond);
		\draw[arr] (cond) -- node[above left, font=\tiny, sloped] {sí} (p2);
		\draw[arr] (p2) -- (p3);
		\draw[arr] (cond) -- node[above right, font=\tiny, sloped] {no} (p3);
		\draw[arr] (p3) -- (p4);
		\draw[arr] (p4) -- (fin);
	\end{tikzpicture}%
	}
	\caption{Diagrama de flujo de las líneas de trabajo futuro propuestas para ViaData — Medellín}
	\label{fig:flujo-trabajo-futuro}
\end{figure}
